\icode{E-MM-MU-L4-261014-002}
光镊是一种利用光力学性质的重要技术。如图所示，真空中放有一个折射率为\(n = \sqrt{3}\)的直角三棱镜，其\(AC\)边长为\(a\)，\(\angle A =\) \,\qty{60}{\degree}，\(\angle B =\) \,\qty{30}{\degree}。一束极细的激光从\(AC\)的中点\(D\)平行于\(BC\)入射。假设激光在\(AC\)面和\(BC\)面不发生反射。已知该束激光的功率为\(P\)。（在真空中的波长为\(\lambda\)，真空中的光速为\(c\)，普朗克常量为\(h\)，\(\sin\) \,\qty{15}{\degree}\( = \frac{\sqrt{6} - \sqrt{2}}{4}\)），则\paren[]
	
\begin{tikzpicture}[decoration={markings,mark=at position 0.5 with {\arrow{>}}},font=\scriptsize]
	\draw (0,0)node[below left]{\(C\)}--(0,2)node[above]{\(A\)}--({2*sqrt(3)},0)node[right]{\(B\)}--(0,0);
	\draw[postaction={decorate},magenta] (-1,1)--(0,1);
	\node at(-.2,.8){\(D\)};
\end{tikzpicture}
\begin{choices}
	\item 激光在棱镜中的传播时间为\(\frac{5a}{2c}\)
	\item 激光在棱镜中的波长为\(\sqrt{3}\lambda\)
	\item 单位时间内通过棱镜的光子数为\(\frac{P\lambda}{hc}\)
	\item 激光对棱镜的作用力大小为\(\frac{(\sqrt{6} - \sqrt{2})P}{2c}\)
\end{choices}
\begin{answer}
    
\end{answer}